\subsection{Постановка задачи}
\label{sec:analysis:specification}

Цель курсового проекта заключается в проектировании и разработке автоматизированной информационной системы управления интернет-магазином автозапчастей. Система обеспечивает ведение каталога деталей, корзину покупателя, транзакционное оформление заказов со списанием остатков, обмен номенклатурой в форматах \textit{JSON}, \textit{XML}, \textit{CSV}, генерацию аналитических отчетов и поддержку мультиязычного интерфейса. База данных реализуется на реляционной СУБД \textit{PostgreSQL} в третьей нормальной форме с организацией доступа через прямые параметризованные SQL-запросы без применения библиотек-посредников. Серверная часть строится на базе веб-фреймворка \textit{FastAPI}. В системе определены две группы пользователей с различными уровнями полномочий: покупатели и администраторы магазина.

В рамках управления каталогом и учетными записями система выполняет регистрацию клиентов, хеширование паролей криптографическим алгоритмом \textit{bcrypt}, авторизацию по электронной почте с выдачей \textit{JWT}-токена и проверку прав доступа к защищенным маршрутам. Для каждой товарной позиции сохраняются исчерпывающие характеристики: наименование, описание, бренд, совместимая модель автомобиля, категория, цена, складской остаток, ссылка на медиаматериалы и статус активности. Обеспечивается постраничная выдача номенклатуры по 20 позиций, сортировка по цене и наименованию, фильтрация по категориям и производителям, и контекстный поиск деталей по артикулу и названию.

Подсистема корзины и оформления заказов реализует добавление товаров с оперативной проверкой доступного количества, динамический пересчет итоговой стоимости, выбор метода оплаты и регистрацию адреса доставки. Процесс покупки выполняется строго в рамках неделимой транзакции СУБД: фиксируются контактные данные клиента, сохраняются контрактные цены позиций на момент совершения сделки и списываются складские запасы. Административный контур обеспечивает ведение ассортимента, модерацию заказов и отслеживание их стадий. При отмене заказа зарезервированные товары автоматически возвращаются на склад. Для синхронизации с внешними платформами поддерживается пакетный импорт и экспорт номенклатуры в форматах \textit{JSON}, \textit{XML} и \textit{CSV} с валидацией структуры данных.

Аналитический модуль системы формирует три вида отчетных документов: товарную накладную по заказу в формате \textit{PDF}, отчет о структуре продаж и доходности позиций в формате \textit{Microsoft Excel} и сводный аналитический отчет для руководства в формате \textit{Microsoft Word}. Документы генерируются потоковым методом в оперативной памяти без сохранения промежуточных файлов на накопителе. Пользовательский веб-интерфейс функционирует в современных браузерах, снабжен визуальной валидацией форм ввода и поддерживает динамическую локализацию на русском и английском языках.

Надежность функционирования обеспечивается соблюдением принципов ACID: сбои на любом этапе оформления заказа приводят к полному откату транзакции без нарушения целостности данных. Ограничения \textit{CHECK} схемы СУБД исключают сохранение отрицательных цен и остатков, а внешние ключи с каскадным удалением предотвращают появление несогласованных записей. Время отклика сервера на операции выборки каталога не превышает 100~мс при объеме номенклатуры до 50\,000 позиций, а транзакция оформления заказа выполняется менее чем за 200~мс за счет построения B-tree индексов по наиболее востребованным полям поиска.

Безопасность достигается обязательной параметризацией всех SQL-запросов для исключения инъекций, надежной изоляцией пользовательских данных и валидацией \textit{JWT}-токенов с ограниченным сроком действия. Для работы клиентской части достаточен персональный компьютер или мобильное устройство с веб-браузером и сетевым подключением. Серверный узел требует наличия процессора с двумя ядрами и тактовой частотой от 2,0~ГГц, оперативной памяти объемом от 4~ГБ и накопителя \textit{SSD} от 20~ГБ. Программный стек сервера включает операционную систему семейства \textit{Linux} или \textit{Windows}, СУБД \textit{PostgreSQL}~15+, интерпретатор \textit{Python}~3.10+, веб-фреймворк \textit{FastAPI} и драйвер \textit{psycopg2-binary}.
